\section{两京失守以后：重建、军镇与区域社会}
\label{sec:late-tang-curated-chronicle}

751年，安禄山兼领范阳、平卢、河东三镇。\eventanchor{ST-751-AN-LUSHAN-COMMANDS}{安禄山兼领三镇（751）}范阳面向幽燕，平卢连着东北，河东控制山西北部的通道；把三镇交给同一指挥中心，意味着将校关系、仓廪与转运机会可被更紧地接合。它不是叛乱的充分原因，却使朝廷在危机中重新分割资源的余地缩小。

755年安禄山起兵，次年两京相继受到冲击。\eventanchor{ST-755-AN-LUSHAN-REBELLION}{安禄山起兵（755）}\eventanchor{LT-756-001}{潼关失守（756）}洛阳、潼关、长安与灵武不是抽象箭头，而是粮秣、关隘、河道和政治名义相互争夺的地点。唐肃宗在灵武立朝、回纥出兵与唐军收复两京，说明唐的重建同时依赖盟友、地方军队和新的财政安排；“平叛”并没有把战前的中央—边镇关系复原。

\begin{turningpoint}{安史之乱改变了什么，而没有改变什么}
危机改变了两京、河北与江淮之间的征调方式，也让节度使成为难以绕开的军政中介。它没有使地方社会只剩一个“藩镇”身份：寺院、商路、流民、州县官与乡里仍在不同资源条件下行动。由战后方镇叙事倒推755年前所有地方都已脱离中央，同样不合材料。
\end{turningpoint}

780年两税法颁行，是战后财政寻找稳定入口的一次选择。\eventanchor{LT-780-020}{两税法（780）}它把征收的语言转向资产和夏秋两期，并不保证每一州县都以同一方式执行。805年永贞革新与819年平淮西，是中枢试图调整宦官、藩镇和行政通道的两次不同尝试。\eventanchor{LT-805-029}{永贞革新（805）}\eventanchor{LT-819-037}{平淮西（819）}它们提示“中央复兴”要经过具体的官员、军队、粮道和地方合作，不能只由年号或皇帝意志判定。

九世纪后半，河北军镇、江淮财源、南诏边事与海路城市的节奏更加分化。874年王仙芝起事、875年黄巢加入起义，将饥荒、流民、地方武装和税收压力推入同一连锁。\eventanchor{LT-874-058}{王仙芝起事（874）}\eventanchor{LT-875-059}{黄巢起事（875）}880年黄巢入长安，唐僖宗西走成都；都城失守不是唐亡的单一原因，却迫使各地军府以自己的粮源和兵员安排应对。\eventanchor{LT-880-062}{黄巢入长安（880）}

\begin{sourcenote}{战乱中的数字、迁徙与地方经验}
两《唐书》《资治通鉴》提供连续年序，却主要从朝廷、将领与政局观察战争；《唐会要》可追踪制度语言，墓志、敦煌文书和城址只给出稀疏的地方切片。兵数、死亡、流民规模和民众动机若无明确材料，不宜写成精确总量。
\end{sourcenote}

\begin{event}{\eventref{LT-780-020}{两税法与战后征收的重新落点}}{建中元年（七八〇年）}{唐与州县、江淮及各地税源}{颁行与制度语言可核；税额、资产认定和地方负担不能由法令推作同一实践}
两税法颁布时，朝廷面对的不是一张等候填入的全国税表，而是战后人口流动、州县财源分化和军费持续索取的局面。把征收安排在夏、秋两期，并以资产等条件重新组织语言，是要让可见的财源进入持续的财政通道；谁来认定资产、如何同旧有负担衔接、军镇和地方中介怎样取用资源，却只能在各地手续中得到不同答案。它为中枢提供新的入口，也把协商、隐匿与摊派的压力带进地方社会。

《旧唐书·德宗纪》、杨炎传、《新唐书》食货材料与《资治通鉴》能核对颁行和朝廷论述；《通典》《唐会要》说明制度术语。它们不能替我们结算一州一县的实收，更不能以一条规范文本证明改革已均匀落地；地方文书和墓志所见的有限切片，正要求保留这种地区差异。
\end{event}

\section{读者事件簇：黄巢余波中的三次改写}

\begin{event}{\eventanchor{LT-882-064}{朱温降唐，宣武军成为新的资源接口}}{中和二年（八八二年）}{唐与宣武 · 黄巢军营、汴州及中原水陆路}{朱温降唐与受授节度使可核；部众规模、授任条件和田令孜的考虑多经后起梁廷叙事塑形}
朱温离开黄巢军而受唐廷授为宣武军节度使时，改变的不只是一个将领的名义。汴州周边的城镇、粮道和运河接点，使他可以把原有的部众、地方征发和朝廷授予的军号编进一套新的供给关系；唐廷也借这一安排寻找能在中原牵制起事军的可用武力。对随军者和沿路居民而言，改换旗号意味着保护、征粮和处罚可能改由不同军府执行，并不等同于战乱已经退去。宣武军因此成为晚唐朝廷与地方兵源重新相接的一处节点。

《旧唐书·朱全忠传》《新五代史·梁本纪第一》及《资治通鉴》唐纪记录了降唐和授任。后梁建立后的回溯性叙事尤其容易把朱温早年的选择写成必然，不能据此确定其全部部众或田令孜的即时盘算。直接后果是唐廷获得一支可调动却不易完全控制的军镇力量；这一不对称关系，后来使朱温能够进入长安并重排宫廷军政。
\end{event}

\begin{event}{\eventanchor{LT-903-077}{朱温入长安，神策军体系被杀戮与改编撕开}}{天复三年（九〇三年）}{唐与朱温、宦官 · 长安宫城、神策军驻地与通往凤翔的道路}{诛宦与军制重组的转折可核；被杀名单、宦官关系及神策军是否完全解体有异文}
朱温率军进入长安后，宫城周围原先由宦官掌握的神策军不再只是朝廷内部的卫队问题。谁能占住城门、营垒和发令处，谁就能决定皇帝身边的护卫、诏令和物资从何处流动；朱温以杀戮和改编切断了宦官原有的接口，也让军人和官员必须在新的军府压力下选择依附或离开。这个行动解决了他进入中枢的即时障碍，却使唐廷更难保有独立于地方强镇的军事支撑。长安因此从皇帝可以调度军队的中心，变成朱温可以借皇帝名义调度资源的场所。

《旧唐书·昭宗本纪》《旧唐书·朱全忠传》与《资治通鉴》唐纪可互校朱温入京、诛宦的顺序。它们对具体名单、宦官与李茂贞的关系及改编程度并不一致，不能把“解体”写作一夜间消失的整套机构。军政接口被改写的后果，却能解释为什么次年迁都洛阳会成为可能：皇帝仍在，保护和命令的路线已经转向。
\end{event}

\begin{event}{\eventanchor{LT-907-081}{后梁建国：开封成为军镇统治的都城}}{开平元年（九〇七年）五月}{后梁 · 开封、汴州水陆转运与河南州县}{建国、改名和定都可核；官制沿革和都城称谓不等于后梁已稳定控制河南以外地区}
朱温改名朱晃、在开封建后梁时，汴州的军镇资源被公开翻译成了都城的名义。开封附近的河道、仓储和道路能把宣武旧部、河南州县与朝廷文书连得更紧，却也使新政权必须持续取得粮饷、船运和地方官的合作；受禅的仪式无法替这些日常成本付款。对仍奉唐号或自保于各镇的人来说，开封的建国并不自动规定他们的选择，反而使中原的多中心竞争更加清楚。八八二年的受授、九〇三年的入京，在这里共同结出一座由军府资源支撑的新都。

《旧五代史·梁太祖纪一》《新五代史·梁本纪第一》与《五代会要》卷一〈帝系〉可核对建国、改名和定都。它们主要从后梁及后出的官修编年观察中枢，不能证明开封以外所有城镇已改从梁廷。后梁建立的直接效果是唐朝名义的终结，中期影响则是军镇、地方财政和旧朝官僚必须在并存政权之间重新选择可以维持道路与供给的关系。
\end{event}

黄巢之后，朱温、李克用及各地军政集团并非只是在唐朝结局中等待出场。它们要面对的税源、城市、河道和边地交涉，将在\eventref{LT-907-080}{后梁建国（907）}以后成为同时展开的多政权历史。

\subsection{江淮盐铁转运：把战后的粮路重新接起来}
\begin{event}{\eventanchor{LT-758-007}{刘晏整顿盐铁转运}}{乾元元年（七五八年）}{唐与江淮转运网络}{受命与整顿方向可核；转运额、盐利与各地负担不能由制度叙事直接换算}
长安收复后，城门重新打开并不等于京师已有粮食。刘晏受命整顿盐铁转运，面对的是江淮仓储、河道、船户、地方官与战后市场如何把粮与财赋送往北方的实际问题。转运网络的价值在于把分散节点接回朝廷，代价却会落在负责装卸、航行、征收和垫付的人身上；它不能被简写成国家财政已经复原。

《旧唐书·刘晏传》《新唐书·刘晏传》与《资治通鉴》可追踪任命和改革的朝廷语言，制度数字多经后出汇整。敦煌、吐鲁番文书以及江淮的地方材料能照见部分行政实践，不能由一处仓或一张文书替整个转运体系结账。
\end{event}